\documentclass[10pt,a4paper]{amsart}
\usepackage{amsmath,amssymb,amsthm,graphicx}
% Retained minimal preamble for the Stacks Project chapter more-algebra.tex (More on Algebra).
% The chapter file pinned-original-source.tex carries only \input{preamble} on its first line, so no
% preamble text lives inside the chapter itself.  The Stacks Project uses one single shared preamble for
% every chapter; the snapshot stages it identically as sources/stacks-chow/preamble.tex and
% sources/stacks-etale-cohomology/preamble.tex.  The block below is copied mechanically, line for line,
% from sources/stacks-chow/preamble.tex lines 14-18 (the Xy-pic setup) and 154-217 (the native theorem
% declarations together with the mathematical macros).  Nothing was added, removed or redefined; the
% wrapper class and the amsmath/amssymb/amsthm/graphicx base packages are supplied by the assembler.
% Chapter: more-algebra.tex -- More on Algebra.
% For commutative diagrams we use Xy-pic
\usepackage[all]{xy}

% We use 2cell for 2-commutative diagrams.
\xyoption{2cell}

% Theorem environments.
%
\theoremstyle{plain}
\newtheorem{theorem}[subsection]{Theorem}
\newtheorem{proposition}[subsection]{Proposition}
\newtheorem{lemma}[subsection]{Lemma}

\theoremstyle{definition}
\newtheorem{definition}[subsection]{Definition}
\newtheorem{example}[subsection]{Example}
\newtheorem{exercise}[subsection]{Exercise}
\newtheorem{situation}[subsection]{Situation}

\theoremstyle{remark}
\newtheorem{remark}[subsection]{Remark}
\newtheorem{remarks}[subsection]{Remarks}

\numberwithin{equation}{subsection}

% Macros
%
\def\lim{\mathop{\mathrm{lim}}\nolimits}
\def\colim{\mathop{\mathrm{colim}}\nolimits}
\def\Spec{\mathop{\mathrm{Spec}}}
\def\Hom{\mathop{\mathrm{Hom}}\nolimits}
\def\Ext{\mathop{\mathrm{Ext}}\nolimits}
\def\SheafHom{\mathop{\mathcal{H}\!\mathit{om}}\nolimits}
\def\SheafExt{\mathop{\mathcal{E}\!\mathit{xt}}\nolimits}
\def\Sch{\mathit{Sch}}
\def\Mor{\mathop{\mathrm{Mor}}\nolimits}
\def\Ob{\mathop{\mathrm{Ob}}\nolimits}
\def\Sh{\mathop{\mathit{Sh}}\nolimits}
\def\NL{\mathop{N\!L}\nolimits}
\def\CH{\mathop{\mathrm{CH}}\nolimits}
\def\proetale{{pro\text{-}\acute{e}tale}}
\def\etale{{\acute{e}tale}}
\def\QCoh{\mathit{QCoh}}
\def\Ker{\mathop{\mathrm{Ker}}}
\def\Im{\mathop{\mathrm{Im}}}
\def\Coker{\mathop{\mathrm{Coker}}}
\def\Coim{\mathop{\mathrm{Coim}}}

% Boxtimes
%
\DeclareMathSymbol{\boxtimes}{\mathbin}{AMSa}{"02}

%
% Macros for moduli stacks/spaces
%
\def\QCohstack{\mathcal{QC}\!\mathit{oh}}
\def\Cohstack{\mathcal{C}\!\mathit{oh}}
\def\Spacesstack{\mathcal{S}\!\mathit{paces}}
\def\Quotfunctor{\mathrm{Quot}}
\def\Hilbfunctor{\mathrm{Hilb}}
\def\Curvesstack{\mathcal{C}\!\mathit{urves}}
\def\Polarizedstack{\mathcal{P}\!\mathit{olarized}}
\def\Complexesstack{\mathcal{C}\!\mathit{omplexes}}
% \Pic is the operator that assigns to X its picard group, usage \Pic(X)
% \Picardstack_{X/B} denotes the Picard stack of X over B
% \Picardfunctor_{X/B} denotes the Picard functor of X over B
\def\Pic{\mathop{\mathrm{Pic}}\nolimits}
\def\Picardstack{\mathcal{P}\!\mathit{ic}}
\def\Picardfunctor{\mathrm{Pic}}
\def\Deformationcategory{\mathcal{D}\!\mathit{ef}}

% --- appended 2026-10-12 (Chief): the project own definitions of the reference/slogan/history
% display environments, copied mechanically from the shared Stacks preamble lines 8-12
% (\usepackage{verbatim} plus the three \newenvironment lines, which the chapter preamble keeps
% verbatim from the project preamble). Nothing added, removed or redefined: these environments
% carry no mathematical content and the project itself typesets them as comments. Without these
% three lines a counted statement whose own text uses them fails with Environment slogan undefined.
\usepackage{verbatim}
\newenvironment{reference}{\comment}{\endcomment}
\newenvironment{slogan}{\comment}{\endcomment}
\newenvironment{history}{\comment}{\endcomment}
\usepackage{cleveref}
\crefname{originallemma}{original lemma}{original lemmas}
\Crefname{originallemma}{Original lemma}{Original lemmas}
\makeatletter
\newlabel{fields-lemma-lift-maps}{{0BME}{source}{}{sourcepointer-1}{}}
\newlabel{fields-lemma-lift-maps@cref}{{[originallemma][0][]0BME}{[1][0][]source}{}{}{}}
\makeatother
\crefname{originallemma}{original lemma}{original lemmas}
\Crefname{originallemma}{Original lemma}{Original lemmas}
\makeatletter
\newlabel{lemma-galois-galois}{{33224}{source}{}{sourcepointer-2}{}}
\newlabel{lemma-galois-galois@cref}{{[originallemma][0][]33224}{[1][0][]source}{}{}{}}
\makeatother
\crefname{originallemma}{original lemma}{original lemmas}
\Crefname{originallemma}{Original lemma}{Original lemmas}
\makeatletter
\newlabel{lemma-one-orbit-geometric-galois-compare}{{32869}{source}{}{sourcepointer-3}{}}
\newlabel{lemma-one-orbit-geometric-galois-compare@cref}{{[originallemma][0][]32869}{[1][0][]source}{}{}{}}
\makeatother
\title{The decomposition, inertia and wild inertia groups of a maximal ideal restrict surjectively along a Galois subtower}
\author{The Stacks Project Authors}
\date{stacks-project:more-algebra.tex:local-snapshot-20261012; GNU FDL 1.2 or later}
\begin{document}
\maketitle
\setcounter{subsection}{1}
\noindent\textit{Editorial note.} All delivered mathematical text is the author's own, copied from the pinned Stacks Project chapter file more-algebra.tex; no statement or proof step was rewritten, reordered or removed.\par\medskip
\noindent\textit{Editorial note.} The chapter file pinned-original-source.tex carries only \textbackslash{}input\{preamble\} on its first line, so no preamble text lives inside the chapter. The retained file sources/stacks-more-algebra/retained-preamble.tex is a mechanical, line-for-line copy of the project's single shared preamble as staged in the snapshot; nothing was added, removed or redefined and no package was added.\par\medskip
\noindent\textit{Editorial note.} The source URL is the source-level Stacks Project URL, as required when more than one card is produced from one source; the per-claim tag is carried by claim\_key.\par\medskip
\noindent\textit{Editorial note.} The counted statement span contains two commutative diagrams and the definition of e' and e; all are delivered verbatim. The three references inside the counted proof span are recorded as non-load-bearing navigation pointers.\par\medskip
\noindent\textit{Allowed standard prerequisites (editorial).}\begin{itemize}
\item Decomposition, inertia and wild inertia groups of a maximal ideal in a Galois extension of a discretely valued field. fixed at statement lines 33624-33632. Are the groups whose restriction behaviour is being described.
\item Normality of the residue field extension and the lifting of automorphisms of a normal subextension. used at proof lines 33663-33668. Produces the right vertical arrow in the first diagram and its commutativity.
\item The inertia character into the roots of unity, and the relation between uniformizers and ramification indices. used at proof lines 33677-33688. Is the computation that gives commutativity of the second diagram up to the power e'/e.
\end{itemize}\medskip
\ifdefined\hypertarget\hypertarget{sourcepointer-1}{}\fi
\noindent\textit{Source reference.} Original lemma 0BME: focused-seeds/AGP/sources/stacks-snapshot-20261011/tags.txt, source lines 14627--14627. Cross-chapter navigation pointer into the Stacks Project; the standard fact it supplies is applied inside the counted proof span and the argument resting on it stays there.\par\medskip
\ifdefined\hypertarget\hypertarget{sourcepointer-2}{}\fi
\noindent\textit{Source reference.} Original lemma 33224: sources/stacks-more-algebra/pinned-original-source.tex, source lines 33224--33224. Same-chapter navigation pointer; the statement it names is standard for the section and the argument that consumes it stays inside the counted proof span.\par\medskip
\ifdefined\hypertarget\hypertarget{sourcepointer-3}{}\fi
\noindent\textit{Source reference.} Original lemma 32869: sources/stacks-more-algebra/pinned-original-source.tex, source lines 32869--32869. Same-chapter navigation pointer; the statement it names is standard for the section and the argument that consumes it stays inside the counted proof span.\par\medskip
\begin{lemma}
\label{lemma-compare-inertia}
Let $A$ be a discrete valuation ring with fraction field $K$.
Let $M/L/K$ be a tower with $M/K$ and $L/K$ finite Galois.
Let $C$, $B$ be the integral closure of $A$ in $M$, $L$.
Let $\mathfrak m' \subset C$ be a maximal ideal and set
$\mathfrak m = \mathfrak m' \cap B$. Let
$$
P \subset I \subset D \subset \text{Gal}(L/K)
\quad\text{and}\quad
P' \subset I' \subset D' \subset \text{Gal}(M/K)
$$
be the wild inertia, inertia, decomposition group of
$\mathfrak m$ and $\mathfrak m'$.
Then the canonical surjection $\text{Gal}(M/K) \to \text{Gal}(L/K)$
induces surjections $P' \to P$, $I' \to I$, and $D' \to D$. Moreover
these fit into commutative diagrams
$$
\vcenter{
\xymatrix{
D' \ar[r] \ar[d] &
\text{Aut}(\kappa(\mathfrak m')/\kappa_A) \ar[d] \\
D \ar[r] &
\text{Aut}(\kappa(\mathfrak m)/\kappa_A)
}
}
\quad\text{and}\quad
\vcenter{
\xymatrix{
I' \ar[r]_-{\theta'} \ar[d] &
\mu_{e'}(\kappa(\mathfrak m')) \ar[d]^{(-)^{e'/e}} \\
I \ar[r]^-\theta &
\mu_e(\kappa(\mathfrak m))
}
}
$$
where $e'$ and $e$ are the ramification indices of
$A \to C_{\mathfrak m'}$ and $A \to B_\mathfrak m$.
\end{lemma}

\begin{proof}
The fact that under the map $\text{Gal}(M/K) \to \text{Gal}(L/K)$
the groups $P', I', D'$ map into $P, I, D$ is immediate from the
definitions of these groups. The commutativity of the first diagram
is clear (observe that since $\kappa(\mathfrak m)/\kappa_A$ is normal
every automorphism of $\kappa(\mathfrak m')$ over $\kappa_A$ indeed
induces an automorphism of $\kappa(\mathfrak m)$ over $\kappa_A$
and hence we obtain the right vertical arrow in the first diagram, see
Lemma \ref{lemma-galois-galois} and
Fields, Lemma \ref{fields-lemma-lift-maps}).

\medskip\noindent
The maps $I' \to I$ and $D' \to D$ are surjective by
Lemma \ref{lemma-one-orbit-geometric-galois-compare}.
The surjectivity of $P' \to P$ follows as $P'$ and $P$
are p-Sylow subgroups of $I'$ and $I$.

\medskip\noindent
To see the commutativity of the second diagram we choose a uniformizer
$\pi'$ of $C_{\mathfrak m'}$ and a uniformizer $\pi$ of $B_\mathfrak m$.
Then $\pi = c' (\pi')^{e'/e}$ for some unit $c'$ of $C_{\mathfrak m'}$.
For $\sigma' \in I'$ the image $\sigma \in I$ is simply the restriction
of $\sigma'$ to $L$. Write $\sigma'(\pi') = c \pi'$ for a unit
$c \in C_{\mathfrak m'}$ and write
$\sigma(\pi) = b \pi$ for a unit $b$ of $B_\mathfrak m$.
Then $\sigma'(\pi) = b \pi$ and we obtain
$$
b \pi = \sigma'(\pi) = \sigma'(c' (\pi')^{e'/e}) =
\sigma'(c') c^{e'/e} (\pi')^{e'/e} =
\frac{\sigma'(c')}{c'} c^{e'/e} \pi
$$
As $\sigma' \in I'$ we see that $b$ and $c^{e'/e}$ have the same
image in the residue field which proves what we want.
\end{proof}

\end{document}
